\section{Introduction}
\IEEEPARstart{C}{ooperative} target protection requires defenders to redistribute their effort as the numbers of vehicles and protected assets change. An attacker approaching a target may already be covered by a well-positioned teammate, while another attacker threatens an exposed asset. Effective actions therefore depend on relationships among attackers, targets, and defenders, interpreted from the perspective of the defender choosing an action. Increasing team size or target count changes these relationships and the distribution of responsibilities. A shared controller that carries useful interaction patterns across engagement sizes would simplify deployment across changing configurations.

Entity-based policies provide a common interface to these changing scenes. Shared encoders and attention accommodate variable entity populations, while REFIL encourages reuse of interaction patterns through randomized entity-wise value factorization~\cite{iqbal2021refil}. Individual decisions also involve two connected computations: developing an entity's context and extracting information relevant to the receiving defender. An attacker's representation can incorporate its relationships to targets and teammates; its relevance to a particular action depends on the defender's own state and history. In the task studied here, public physical states describe the scene, while the attacker's private target assignment gives recent motion a role in interpreting that scene.

Shared iterative attention offers a reusable operator for developing entity context. Relational reinforcement learning and Universal Transformers establish repeated processing with shared attention parameters~\cite{zambaldi2019relational,dehghani2019universal}. Each fully connected update accesses the permitted entity set, and subsequent updates process the context formed by earlier applications. Preserving intermediate states gives the decision stage access to several stages of this computation, following the broader principle of combining representations across depth~\cite{xu2018jk}. Figure~\ref{fig:motivation} illustrates the resulting design idea: develop entity context and make its intermediate representations available to an individualized reader.

We introduce \textbf{Looped Entity Attention Fusion (LEAF)}, which connects shared entity refinement to individual action selection through history-conditioned reading and learned cross-round fusion. Each defender initializes a receiver-relative entity workspace and repeatedly applies one shared attention block. A query constructed from its initial self representation and previous temporal memory reads every intermediate entity set. Learned weights combine these summaries, and a residual projection adds the fused information to the feature entering a recurrent utility network. The network then performs one temporal update and produces action values. Figure~\ref{fig:framework} shows how within-step refinement connects to the memory carried between physical decisions.

Training samples the unfolding depth while retaining REFIL's entity-factorized cooperative learning objective. Early representations directly supply the reader and also become inputs to subsequent updates at longer unfoldings. The same policy parameters therefore learn through both short decision paths and longer sequences of entity processing. Training across depths has a related precedent in LayerDrop~\cite{fan2020layerdrop}; LEAF uses different unfolding lengths of a shared operator with access to the resulting intermediate states. This design motivates evaluating the quality of a learned one-round policy alongside the effects of additional execution rounds and different readout choices.

We evaluate frozen-policy transfer across team-size, target-count, and joint changes in simulated protection against a fixed reactive attacker controller (Fig.~\ref{fig:scenario}). LEAF improves protection relative to REFIL across the evaluated population shifts. Matched one-round execution supports a more effective policy after multi-depth training, while frozen-policy interventions connect the learned decisions to early reads and temporal conditioning. Representation probes further examine the geometric information accessible in the entity states. The contributions are:
\begin{enumerate}
\item A shared refinement and individual readout architecture that connects intermediate entity states to recurrent action values.
\item A sampled-depth learning design that couples short and long unfoldings within entity-factorized value learning.
\item A population-transfer evaluation with matched-depth comparisons and targeted interventions on intermediate reads and temporal conditioning.
\end{enumerate}
